% GENERATED FILE. Edit canonical lessons, then run npm run generate:books.
% canonical-source-hash: fnv1a64:dc6e21d0e160185b
% canonical-lessons: TE-C267-indradhanussu, TE-C267-varada, TE-C267-karuvu, TE-C267-vennela, TE-C267-cikati

\chapter{Rainbow, Flood, Drought, Moonlight, Darkness}
\label{ch:te-rainbow-flood-drought-moonlight-darkness}
\hlchaptermodality{267}

\begin{chapteropening}
\textbf{By the end of this chapter:} \emph{I can say a rainbow, a flood, a drought, moonlight and darkness.}

Complete the last lesson of chapter 267: I can say a rainbow, a flood, a drought, moonlight and darkness.
\end{chapteropening}

\section[\texorpdfstring{\te{ఇంద్రధనుస్సు} (indradhanussu)}{ఇంద్రధనుస్సు (indradhanussu)}]{\te{ఇంద్రధనుస్సు} (indradhanussu) — a rainbow}

\label{lesson:TE-C267-indradhanussu}



\begin{quote}
Before the new one: say the Telugu for holy basil, then the Telugu for a bud.
\end{quote}

\subsection*{You'll want to know: \te{ఇంద్రధనుస్సు}}
\textbf{\te{ఇంద్రధనుస్సు}} — \emph{indradhanussu} — \textquotedblleft{}a rainbow\textquotedblright{}.

\begin{grammarlens}[title={What you've built}]
One more word for this chapter.
\end{grammarlens}

\subsection*{Your turn}
\begin{itemize}
\raggedright
  \item \emph{Say it:} \emph{indradhanussu}
  \item \emph{Say it:} \emph{indradhanussu}, once more
  \item \emph{Say it:} say \emph{indradhanussu}
  \item \emph{Recall:} say the Telugu for holy basil, then the Telugu for a bud, then say \emph{indradhanussu} again
\end{itemize}

\begin{tcolorbox}[breakable,colback=teal!4,colframe=teal!35!black,title={Before you move on}]
What does \te{ఇంద్రధనుస్సు} mean? (\textquotedblleft{}A rainbow\textquotedblright{}.) Say it once more.
\end{tcolorbox}
\section[\texorpdfstring{\te{వరద} (varada)}{వరద (varada)}]{\te{వరద} (varada) — a flood}

\label{lesson:TE-C267-varada}



\begin{quote}
Before the new one: say the Telugu for a bud, then the Telugu for a rainbow.
\end{quote}

\subsection*{You'll want to know: \te{వరద}}
\textbf{\te{వరద}} — \emph{varada} — \textquotedblleft{}a flood\textquotedblright{}.

\begin{grammarlens}[title={What you've built}]
One more word for this chapter.
\end{grammarlens}

\subsection*{Your turn}
\begin{itemize}
\raggedright
  \item \emph{Say it:} \emph{varada}
  \item \emph{Say it:} \emph{varada}, once more
  \item \emph{Say it:} say \emph{varada}
  \item \emph{Recall:} say the Telugu for a bud, then the Telugu for a rainbow, then say \emph{varada} again
\end{itemize}

\begin{tcolorbox}[breakable,colback=teal!4,colframe=teal!35!black,title={Before you move on}]
What does \te{వరద} mean? (\textquotedblleft{}A flood\textquotedblright{}.) Say it once more.
\end{tcolorbox}
\section[\texorpdfstring{\te{కరువు} (karuvu)}{కరువు (karuvu)}]{\te{కరువు} (karuvu) — a drought, a famine}

\label{lesson:TE-C267-karuvu}



\begin{quote}
Before the new one: say the Telugu for a rainbow, then the Telugu for a flood.
\end{quote}

\subsection*{You'll want to know: \te{కరువు}}
\textbf{\te{కరువు}} — \emph{karuvu} — \textquotedblleft{}a drought, a famine\textquotedblright{}.

\begin{grammarlens}[title={What you've built}]
One more word for this chapter.
\end{grammarlens}

\subsection*{Your turn}
\begin{itemize}
\raggedright
  \item \emph{Say it:} \emph{karuvu}
  \item \emph{Say it:} \emph{karuvu}, once more
  \item \emph{Say it:} say \emph{karuvu}
  \item \emph{Recall:} say the Telugu for a rainbow, then the Telugu for a flood, then say \emph{karuvu} again
\end{itemize}

\begin{tcolorbox}[breakable,colback=teal!4,colframe=teal!35!black,title={Before you move on}]
What does \te{కరువు} mean? (\textquotedblleft{}A drought, a famine\textquotedblright{}.) Say it once more.
\end{tcolorbox}
\section[\texorpdfstring{\te{వెన్నెల} (vennela)}{వెన్నెల (vennela)}]{\te{వెన్నెల} (vennela) — moonlight}

\label{lesson:TE-C267-vennela}



\begin{quote}
Before the new one: say the Telugu for a flood, then the Telugu for a drought.
\end{quote}

\subsection*{You'll want to know: \te{వెన్నెల}}
\textbf{\te{వెన్నెల}} — \emph{vennela} — \textquotedblleft{}moonlight\textquotedblright{}.

\begin{grammarlens}[title={What you've built}]
One more word for this chapter.
\end{grammarlens}

\subsection*{Your turn}
\begin{itemize}
\raggedright
  \item \emph{Say it:} \emph{vennela}
  \item \emph{Say it:} \emph{vennela}, once more
  \item \emph{Say it:} say \emph{vennela}
  \item \emph{Recall:} say the Telugu for a flood, then the Telugu for a drought, then say \emph{vennela} again
\end{itemize}

\begin{tcolorbox}[breakable,colback=teal!4,colframe=teal!35!black,title={Before you move on}]
What does \te{వెన్నెల} mean? (\textquotedblleft{}Moonlight\textquotedblright{}.) Say it once more.
\end{tcolorbox}
\section[\texorpdfstring{\te{చీకటి} (cīkaṭi)}{చీకటి (cīkaṭi)}]{\te{చీకటి} (cīkaṭi) — darkness}

\label{lesson:TE-C267-cikati}



\begin{quote}
Before the new one: say the Telugu for a drought, then the Telugu for moonlight.
\end{quote}

\subsection*{You'll want to know: \te{చీకటి}}
\textbf{\te{చీకటి}} — \emph{cīkaṭi} — \textquotedblleft{}darkness\textquotedblright{}.

\begin{grammarlens}[title={What you've built}]
One more word for this chapter.
\end{grammarlens}

\subsection*{Your turn}
\begin{itemize}
\raggedright
  \item \emph{Say it:} \emph{cīkaṭi}
  \item \emph{Say it:} \emph{cīkaṭi}, once more
  \item \emph{Say it:} say \emph{cīkaṭi}
  \item \emph{Recall:} say the Telugu for a drought, then the Telugu for moonlight, then say \emph{cīkaṭi} again
\end{itemize}

\begin{tcolorbox}[breakable,colback=teal!4,colframe=teal!35!black,title={Before you move on}]
What does \te{చీకటి} mean? (\textquotedblleft{}Darkness\textquotedblright{}.) Say it once more.
\end{tcolorbox}
